Mapa daňového klínu na práci v~\ac{EU}27 (Eurostat, strukturální ukazatel pro svobodného bezdětného zaměstnance při \SI{67}{\percent} průměrné mzdy) ukazuje, že ČR se nachází mírně nad mediánem --- bez ohledu na to, že nominální sazby daně z~příjmu jsou relativně nízké, protože celkový klín zahrnuje zaměstnanecké i~zaměstnavatelské odvody.
Tento relativně vysoký klín v~kontextu nízké mzdy vytváří specifické napětí: zaměstnavatel platí za hodinu práce relativně vysoké odvody, ale pracovník dostává relativně málo --- mezera jde na vrub odvodu, nikoli na přidanou hodnotu pro pracovníka.
Státy s~vysokým pokrytím \ac{KS} jako AT a~DK sice vykazují srovnatelný nebo vyšší klín, ale kompenzují jej vyšší hrubou mzdou sjednanou kolektivní smlouvou: de facto funguje model „vyšší odvody za vyšší záruku mzdy a~kolektivní ochrany“.
V~ČR je klín podobný, ale hrubá mzda nižší --- pracovník tedy nesetí výhody ani nízkého klínu (jako třeba SK), ani vysoké hrubé mzdy (jako AT, DK): argument pro reformu kolektivního vyjednávání s~cílem zvýšit hrubou mzdu a~ospravedlnit výši klínu přínosy pro pracovníka.

Mapa daňového klínu na práci v~\ac{EU}27 (Eurostat, strukturální ukazatel pro svobodného bezdětného zaměstnance při \SI{67}{\percent} průměrné mzdy) ukazuje, že ČR se nachází mírně nad mediánem --- bez ohledu na to, že nominální sazby daně z~příjmu jsou relativně nízké, protože celkový klín zahrnuje zaměstnanecké i~zaměstnavatelské odvody.
Tento relativně vysoký klín v~kontextu nízké mzdy vytváří specifické napětí: zaměstnavatel platí za hodinu práce relativně vysoké odvody, ale pracovník dostává relativně málo --- mezera jde na vrub odvodu, nikoli na přidanou hodnotu pro pracovníka.
Státy s~vysokým pokrytím \ac{KS} jako AT a~DK sice vykazují srovnatelný nebo vyšší klín, ale kompenzují jej vyšší hrubou mzdou sjednanou kolektivní smlouvou: de facto funguje model „vyšší odvody za vyšší záruku mzdy a~kolektivní ochrany“.
V~ČR je klín podobný, ale hrubá mzda nižší --- pracovník tedy nesetí výhody ani nízkého klínu (jako třeba SK), ani vysoké hrubé mzdy (jako AT, DK): argument pro reformu kolektivního vyjednávání s~cílem zvýšit hrubou mzdu a~ospravedlnit výši klínu přínosy pro pracovníka.

Mapa daňového klínu na práci v~\ac{EU}27 (Eurostat, strukturální ukazatel pro svobodného bezdětného zaměstnance při \SI{67}{\percent} průměrné mzdy) ukazuje, že ČR se nachází mírně nad mediánem --- bez ohledu na to, že nominální sazby daně z~příjmu jsou relativně nízké, protože celkový klín zahrnuje zaměstnanecké i~zaměstnavatelské odvody.
Tento relativně vysoký klín v~kontextu nízké mzdy vytváří specifické napětí: zaměstnavatel platí za hodinu práce relativně vysoké odvody, ale pracovník dostává relativně málo --- mezera jde na vrub odvodu, nikoli na přidanou hodnotu pro pracovníka.
Státy s~vysokým pokrytím \ac{KS} jako AT a~DK sice vykazují srovnatelný nebo vyšší klín, ale kompenzují jej vyšší hrubou mzdou sjednanou kolektivní smlouvou: de facto funguje model „vyšší odvody za vyšší záruku mzdy a~kolektivní ochrany“.
V~ČR je klín podobný, ale hrubá mzda nižší --- pracovník tedy nesetí výhody ani nízkého klínu (jako třeba SK), ani vysoké hrubé mzdy (jako AT, DK): argument pro reformu kolektivního vyjednávání s~cílem zvýšit hrubou mzdu a~ospravedlnit výši klínu přínosy pro pracovníka.

Mapa daňového klínu na práci v~\ac{EU}27 (Eurostat, strukturální ukazatel pro svobodného bezdětného zaměstnance při \SI{67}{\percent} průměrné mzdy) ukazuje, že ČR se nachází mírně nad mediánem --- bez ohledu na to, že nominální sazby daně z~příjmu jsou relativně nízké, protože celkový klín zahrnuje zaměstnanecké i~zaměstnavatelské odvody.
Tento relativně vysoký klín v~kontextu nízké mzdy vytváří specifické napětí: zaměstnavatel platí za hodinu práce relativně vysoké odvody, ale pracovník dostává relativně málo --- mezera jde na vrub odvodu, nikoli na přidanou hodnotu pro pracovníka.
Státy s~vysokým pokrytím \ac{KS} jako AT a~DK sice vykazují srovnatelný nebo vyšší klín, ale kompenzují jej vyšší hrubou mzdou sjednanou kolektivní smlouvou: de facto funguje model „vyšší odvody za vyšší záruku mzdy a~kolektivní ochrany“.
V~ČR je klín podobný, ale hrubá mzda nižší --- pracovník tedy nesetí výhody ani nízkého klínu (jako třeba SK), ani vysoké hrubé mzdy (jako AT, DK): argument pro reformu kolektivního vyjednávání s~cílem zvýšit hrubou mzdu a~ospravedlnit výši klínu přínosy pro pracovníka.

\input{texparts/python/tax_wedge_map}



